\begin{EnChapter}

\chapter{System Model}
\label{chap:system-model}

This chapter develops a topology-agnostic downlink model of an O-RAN-controlled
integrated access and backhaul (IAB) network. Section~\ref{subsec:topology}
represents a multi-hop IAB network of arbitrary depth as a rooted tree.
Sections~\ref{subsec:resource}--\ref{subsec:coupling} describe the radio resources,
the link rates, and the scheduling of the distributed unit (DU) of each node; each
IAB node also contains a mobile termination (MT) that connects it to its parent. Section~\ref{subsec:traffic} models multi-hop
traffic and queues, and Section~\ref{subsec:problem} states the optimization
problem. Section~\ref{subsec:marl} describes the multi-agent reinforcement learning
(MARL) formulation used to approximate its solution, Section~\ref{subsec:fl}
describes the hierarchical federated learning (FL) built on top of it, and
Section~\ref{subsec:oran} maps each component to the O-RAN architecture.

\section{Network Topology}
\label{subsec:topology}

Consider a downlink IAB network consisting of a single IAB donor, $N$ IAB nodes
and $U$ user equipments (UEs). Let
\begin{equation}
  \Nset = \{n_1, \dots, n_N\}, \qquad \Uset = \{u_1, \dots, u_U\}
\end{equation}
be the sets of IAB nodes and UEs, where all $N+U$ elements are pairwise distinct,
and let $0$ denote the donor, with $0 \notin \Nset \cup \Uset$. The network is
represented by a rooted directed tree
\begin{equation}
  \mathbb{T} = (\Vset, \Eset), \qquad \Vset = \{0\} \cup \Nset \cup \Uset,
\end{equation}
where the union is disjoint. Throughout, $n,m$ denote generic elements of
$\{0\} \cup \Nset$, $u$ a generic element of $\Uset$, and $v$ a generic element of
$\Nset \cup \Uset$.

\begin{definition}[Parent and children]
\label{def:parent}
The parent map $\parent : \Nset \cup \Uset \rightarrow \{0\} \cup \Nset$ assigns
each IAB node and each UE to its serving node, i.e., the node whose DU the MT of the
IAB node, or the UE, is attached to. The edge set is
$\Eset = \{(\parent(v), v) : v \in \Nset \cup \Uset\}$, and each edge is a downlink
radio link served by the DU of $\parent(v)$. For $n \in \{0\} \cup \Nset$, the set of
children is
\begin{equation}
  \Cset_n = \Cset_n^{\mathrm{I}} \cup \Cset_n^{\mathrm{U}}, \quad
  \Cset_n^{\mathrm{I}} = \{m \in \Nset : \parent(m) = n\}, \quad
  \Cset_n^{\mathrm{U}} = \{u \in \Uset : \parent(u) = n\}.
\end{equation}
\end{definition}

The model imposes no restriction on the shape of the topology. Any node may have
both IAB children and UE children ($\Cset_n^{\mathrm{I}} \neq \emptyset$ and
$\Cset_n^{\mathrm{U}} \neq \emptyset$), and both the fan-out and the depth are
arbitrary. The only structural requirements are that $\mathbb{T}$ is a tree rooted
at $0$, so that every node has exactly one path to the donor (the single-parent
case of IAB), and that UEs are leaves.

\paragraph{Depth, paths and subtrees.}
The depth of a vertex is $h(0)=0$ and $h(v) = h(\parent(v)) + 1$. The maximum depth
$H = \max_{u \in \Uset} h(u)$ is unrestricted, and the hop count of UE $u$ is
$h(u)$. Let $\parent^{(j)}(\cdot)$ denote the $j$-th ancestor, with
$\parent^{(0)}(v) = v$. The downlink path of UE $u$ is
\begin{equation}
  \mathcal{P}(u) = \bigl\{ (\parent^{(j+1)}(u), \parent^{(j)}(u)) : j = 0, \dots, h(u)-1 \bigr\}.
\end{equation}
Let $\Uset_v \subseteq \Uset$ be the UEs in the subtree rooted at $v$, with
$\Uset_u = \{u\}$ for $u \in \Uset$, and let $\Nset_v$ be the IAB nodes in that
subtree, including $v$ itself when $v \in \Nset$. The link $(\parent(v), v)$ carries
exactly the downlink traffic destined to $\Uset_v$.

\paragraph{Roles and branches.}
Two topological attributes, used by the federated aggregation in
Section~\ref{subsec:fl}, are defined from a node's position in the tree. For
$n \in \Nset$,
\begin{align}
  \role(n) &=
  \begin{cases}
    \textsf{relay},  & \Cset_n^{\mathrm{I}} \neq \emptyset, \\
    \textsf{access}, & \Cset_n^{\mathrm{I}} = \emptyset,
  \end{cases}
  \label{eq:role} \\
  \br(n) &= \parent^{(h(n)-1)}(n) \in \Cset_0^{\mathrm{I}}.
  \label{eq:branch}
\end{align}
A node is a relay if and only if it has downstream IAB nodes, independently of its
depth, and a relay may also serve UEs directly. $\br(n)$ is the ancestor of $n$ at
depth one. The set of branches is $\Bset = \Cset_0^{\mathrm{I}}$, and branch
$b \in \Bset$ contains the IAB nodes $\Nset_b$.

\paragraph{Node architecture.}
Each IAB node $n \in \Nset$ consists of an MT and a DU. MT$_n$ attaches to
DU$_{\parent(n)}$ as a regular UE to obtain backhaul, and DU$_n$ serves every child
in $\Cset_n$. The donor consists of a central unit (CU) and DU$_0$ and is connected
to the core network through a wired interface. From the viewpoint of the scheduler
of DU$_n$, every child $v \in \Cset_n$ is one scheduling entity, whether it is a UE
or the MT of a downstream IAB node.

\begin{assumption}[Backhaul forwarding and queue structure]
\label{ass:forwarding}
The Backhaul Adaptation Protocol (BAP) defined for standard IAB in 3GPP TS 38.340~\cite{3gpp38340} is
not used. Instead, each MT establishes a single protocol data unit (PDU) session
with a single data radio bearer, and traffic is forwarded between hops by IP
routing over that bearer. Consequently DU$_n$ maintains exactly one radio link
control (RLC) queue per child $v \in \Cset_n$, and for an IAB child
$m \in \Cset_n^{\mathrm{I}}$ this queue aggregates the traffic of all UEs in
$\Uset_m$ in first-in-first-out order.
\end{assumption}

\begin{example}[Testbed topology]
\label{ex:testbed}
The testbed is the instance $N=12$, $U=16$, $H=3$ with
\begin{itemize}
  \item $\Cset_0^{\mathrm{I}} = \{n_1, n_2, n_3, n_4\}$ and
        $\Cset_0^{\mathrm{U}} = \emptyset$;
  \item for $i=1,\dots,4$: $\Cset_{n_i}^{\mathrm{I}} = \{n_{2i+3}, n_{2i+4}\}$ and
        $\Cset_{n_i}^{\mathrm{U}} = \emptyset$, so $\role(n_i)=\textsf{relay}$;
  \item for $i=5,\dots,12$: $\Cset_{n_i}^{\mathrm{I}} = \emptyset$ and
        $\Cset_{n_i}^{\mathrm{U}} = \{u_{2i-9}, u_{2i-8}\}$, so
        $\role(n_i)=\textsf{access}$;
  \item $|\Bset| = 4$ and $\Nset_{n_i} = \{n_i, n_{2i+3}, n_{2i+4}\}$ for
        $i=1,\dots,4$.
\end{itemize}
All UEs are three hops away from the donor. Other configurations allowed by the
general model, such as a relay serving UEs directly, branches of unequal depth, or
$H>3$, do not occur in this instance; the equations below apply to them without
modification.
\end{example}

\section{Time Structure and Radio Resources}
\label{subsec:resource}

Time is slotted with index $t \in \{0,1,2,\dots\}$ and slot length
$T_{\mathrm{slot}} = 2^{-\mu}$~ms, where $\mu$ is the numerology. Each DU$_n$ owns
$B_n$ physical resource blocks (PRBs) of 12 subcarriers on its carrier. All carriers
use time-division duplexing (TDD) with a common periodic pattern of $P$ slots, of
which $P^{\mathrm{DL}}$ slots can carry downlink data and $P^{\mathrm{UL}}$ slots
can carry uplink data. Let $\Tset^{\mathrm{DL}}$ denote the set of downlink-capable
slots and
\begin{equation}
  \delta^{\mathrm{DL}} = \frac{P^{\mathrm{DL}}}{P}, \qquad
  \delta^{\mathrm{UL}} = \frac{P^{\mathrm{UL}}}{P}
\end{equation}
the corresponding fractions of time. Control decisions are made per epoch. Epoch
$k$ is the interval $[t_k, t_{k+1})$ of length $\Delta$ seconds and contains the
slot set $\Tset_k$.

\begin{assumption}[Out-of-band backhaul and orthogonal carriers]
\label{ass:orthogonal}
The backhaul and access links use separate frequency resources, i.e., the IAB nodes
operate in out-of-band mode~\cite{topcu2025oranIAB,moro2023openIAB}, so the MT
and the DU of a node can receive and transmit simultaneously. In addition, the DUs operate
on non-overlapping carriers, so there is no co-channel interference between links.
This lets the model focus on traffic aggregation and scheduling over multi-hop
backhaul. In-band operation, in which the MT and the DU of a node share the same
resources under a half-duplex constraint, is outside the scope of this work; it
would enter the model through the available PRBs in~\eqref{eq:beta} and through
co-channel links in the interference term $I_v(t)$ in~\eqref{eq:sinr}.
\end{assumption}

\section{Channel and Link Rate}
\label{subsec:rate}

The per-PRB signal-to-interference-plus-noise ratio (SINR) of link
$(\parent(v), v)$ in slot $t$ is
\begin{equation}
  \gamma_v(t) = \frac{P_{\parent(v)}\, g_v(t)}{I_v(t) + \sigma^2},
  \label{eq:sinr}
\end{equation}
where $P_{\parent(v)}$ is the per-PRB transmit power, $g_v(t)$ contains path loss
and small-scale fading, $I_v(t)$ is the interference ($I_v \equiv 0$ under
Assumption~\ref{ass:orthogonal}), and $\sigma^2$ is the noise power. The DU performs
link adaptation on the reported channel quality and selects a modulation and coding
scheme (MCS) index $q_v(t)$, which maps to a modulation order $Q_m(q)$ and a code
rate $R(q)$ according to the MCS tables of 3GPP TS 38.214~\cite{3gpp38214}. The number of bits
carried by one PRB in one downlink slot is
\begin{equation}
  \rho_v(t) = \nu_v\, Q_m\bigl(q_v(t)\bigr)\, R\bigl(q_v(t)\bigr)\, N_{\mathrm{sc}}^{\mathrm{RB}}\, N_{\mathrm{symb}}\, (1 - \mathrm{OH}),
  \label{eq:rho}
\end{equation}
where $\nu_v$ is the number of MIMO layers, $N_{\mathrm{sc}}^{\mathrm{RB}}=12$,
$N_{\mathrm{symb}}=14$ is the number of OFDM symbols per slot, and $\mathrm{OH}$ is
the control and reference-signal overhead. If DU$_{\parent(v)}$ allocates
$b_{\parent(v),v}(t)$ PRBs to $v$ in slot $t$, at most $b_{\parent(v),v}(t)\,\rho_v(t)$
bits are delivered on the link in that slot.

\paragraph{Peak rate.}
Evaluating~\eqref{eq:rho} at the highest MCS on every symbol gives the approximate
maximum data rate of 3GPP TS 38.306~\cite{3gpp38306},
\begin{equation}
  r^{\max}_n = 10^{-6}\, \nu\, Q_m^{\max}\, f\, R^{\max}\, \frac{12\, B_n}{T_{\mathrm{s}}^{\mu}}\, (1 - \mathrm{OH})
  \ \ \text{[Mbps]}, \qquad
  T_{\mathrm{s}}^{\mu} = \frac{10^{-3}}{14 \cdot 2^{\mu}}\ \text{s},
  \label{eq:peak}
\end{equation}
where $f$ is the scaling factor. This expression assumes that every symbol carries
downlink data; under TDD the downlink peak is approximately
$\delta^{\mathrm{DL}}_{\mathrm{symb}}\, r^{\max}_n$, where
$\delta^{\mathrm{DL}}_{\mathrm{symb}}$ is the fraction of downlink symbols. Both are
upper bounds on the MAC-layer rate and are not attainable application-layer
throughputs.

\section{DU Scheduling}
\label{subsec:coupling}

\paragraph{DU resource constraint.}
The numbers of PRBs $b_{n,v}(t) \in \mathbb{Z}_{\ge 0}$ that DU$_n$ allocates to its
children in a downlink slot $t \in \Tset^{\mathrm{DL}}$ satisfy
\begin{equation}
  \sum_{v \in \Cset_n} b_{n,v}(t) \le \tilde B_n(t), \qquad \forall n \in \{0\} \cup \Nset,
  \label{eq:du-budget}
\end{equation}
where $\tilde B_n(t) \le B_n$ is the number of PRBs available to DU$_n$ in that
slot. For slots that cannot carry downlink data, $t \notin \Tset^{\mathrm{DL}}$, we
set $b_{n,v}(t) = 0$ and $\tilde B_n(t) = 0$, so that all time averages below are
taken over all slots. In downlink-capable slots,
\begin{equation}
  \tilde B_n(t) = \bigl\lfloor \beta_n(t)\, B_n \bigr\rfloor, \qquad
  \beta_n(t) \in [0,1], \qquad t \in \Tset^{\mathrm{DL}},
  \label{eq:beta}
\end{equation}
where $\beta_n(t)$ is the fraction of the PRBs of DU$_n$ that the resource
configuration makes available. In an in-band IAB node, $\beta_n(t)$ would be set by
the Hard/Soft/Not-Available (H/S/NA) configuration that keeps the DU off the
resources used by its MT. Under Assumption~\ref{ass:orthogonal}, the MT and the DU
do not share resources, so $\beta_n(t) \equiv 1$ and $\tilde B_n(t) = B_n$ for all
$n \in \{0\} \cup \Nset$. The general form~\eqref{eq:beta} is kept so that the
observation of Section~\ref{subsec:marl} is defined for both cases.

\paragraph{Scheduling policy.}
Within the budget~\eqref{eq:du-budget}, DU$_n$ schedules with proportional fair (PF)
in every downlink slot. Child $v$ has priority $M_v(t) = \rho_v(t) / \bar R_v(t)$,
where $\bar R_v$ is an exponentially smoothed average of its served rate. The
scheduler allocates in decreasing priority, with at least $b_{\min}$ PRBs per
scheduled entity, until $\tilde B_n(t)$ is exhausted or all queues are empty. The
control variable of this work is a per-child PRB cap imposed on PF,
\begin{equation}
  b_{n,v}(t) \le \bar b_{n,v}(k) =
  \begin{cases}
    \max\bigl\{ \lfloor a_{n,v}(k)\, B_n \rfloor,\ 1 \bigr\}, & a_{n,v}(k) < 1, \\
    B_n, & a_{n,v}(k) = 1,
  \end{cases}
  \qquad \forall t \in \Tset_k,
  \label{eq:cap}
\end{equation}
where $a_{n,v}(k) \in \Aset = \{\alpha_1 < \alpha_2 < \dots < \alpha_L = 1\}$ is a
discrete tier. The cap is held constant during epoch $k$, and PF still allocates
within the caps. When $a_{n,v} = 1$ for all $v \in \Cset_n$, the caps are inactive
and DU$_n$ behaves exactly as PF, so PF is a particular policy of the control space.
A cap is active in slot $t$ when the allocation that PF would give to child $v$
without the cap exceeds $\bar b_{n,v}(k)$; this can happen even when the total demand
of the children is below $\tilde B_n(t)$. An active cap reduces the throughput of $v$
if the demand of $v$ exceeds the rate supported by the cap, and otherwise only
spreads its transmissions over more slots and increases its delay. Because $\Cset_n$ contains both UEs and downstream IAB nodes, the
action of a relay also decides how its resources are split between directly
attached UEs and backhaul links.

\section{Multi-Hop Traffic and Queues}
\label{subsec:traffic}

\paragraph{Sources.}
Each UE $u$ is associated with a downlink flow from the core network with a
target (offered) rate $\lambda_u(k) \ge 0$ in epoch $k$. The bits injected into
the donor in slot $t$, $A_u(t)$, are determined by the transport protocol. For a
rate-limited TCP flow, the sender injects at most $\lambda_u(k)$ on average, and
congestion control reduces the injection when end-to-end delivery is slower, so
excess demand is not admitted into the network. For a UDP flow, $A_u(t)$ follows
the target rate regardless of delivery, and excess traffic is dropped at finite
buffers.

\paragraph{Queue dynamics.}
For each node $n \in \{0\} \cup \Nset$ and UE $u \in \Uset_n$, let $Q_n^{(u)}(t)$ be
the bits of $u$ buffered at DU$_n$ and waiting for the next hop
$\mathrm{nh}_n(u) \in \Cset_n$, i.e., the child of $n$ on $\mathcal{P}(u)$. Let
$S_n^{(u)}(t)$ be the bits of $u$ sent by DU$_n$ in slot $t$ and $A_n^{(u)}(t)$ the
bits of $u$ arriving at DU$_n$, with
\begin{equation}
  A_0^{(u)}(t) = A_u(t), \qquad
  A_n^{(u)}(t) = S_{\parent(n)}^{(u)}(t - \tau_n), \quad n \in \Nset,
  \label{eq:flow-conservation}
\end{equation}
where $\tau_n$ is the forwarding delay. Equation~\eqref{eq:flow-conservation} is
per-hop flow conservation. By Assumption~\ref{ass:forwarding}, DU$_n$ keeps one
finite buffer of capacity $Q^{\mathrm{buf}}$ per child $v \in \Cset_n$, holding the
aggregated queue $Q_{n,v}(t) = \sum_{u \in \Uset_v} Q_n^{(u)}(t)$ with aggregated
arrivals $A_{n,v}(t) = \sum_{u \in \Uset_v} A_n^{(u)}(t)$. The service, the buffer
update and the dropped bits are
\begin{align}
  S_{n,v}(t) &= \min\bigl\{ b_{n,v}(t)\,\rho_v(t),\ Q_{n,v}(t) \bigr\}, \label{eq:queue-agg}\\
  Q_{n,v}(t+1) &= \min\Bigl\{ Q^{\mathrm{buf}},\ \bigl[Q_{n,v}(t) - S_{n,v}(t)\bigr]^{+} + A_{n,v}(t) \Bigr\}, \label{eq:queue}\\
  D_{n,v}(t) &= \Bigl[ \bigl[Q_{n,v}(t) - S_{n,v}(t)\bigr]^{+} + A_{n,v}(t) - Q^{\mathrm{buf}} \Bigr]^{+}, \label{eq:drop}
\end{align}
where $[\cdot]^{+} = \max\{\cdot, 0\}$; the inner projection is redundant
under~\eqref{eq:queue-agg} but makes the recursion well defined for any service
rule. Hence $0 \le Q_{n,v}(t) \le Q^{\mathrm{buf}}$ for all $t$. The buffer is served
in first-in-first-out order and overflow is handled by tail drop: $S_{n,v}(t)$ is
split among $\Uset_v$ in arrival order, and the dropped bits $D_{n,v}(t)$ are the
most recently arrived ones. This defines the per-UE service $S_n^{(u)}(t)$ and drops
$D_n^{(u)}(t)$, with $\sum_{u \in \Uset_v} D_n^{(u)}(t) = D_{n,v}(t)$, and
\begin{equation}
  Q_n^{(u)}(t+1) = Q_n^{(u)}(t) - S_n^{(u)}(t) + A_n^{(u)}(t) - D_n^{(u)}(t).
\end{equation}
The throughput of UE $u$ in epoch $k$ is what the last hop
delivers,
\begin{equation}
  x_u(k) = \frac{1}{\Delta} \sum_{t \in \Tset_k} S_{\parent(u),u}(t).
  \label{eq:ue-tput}
\end{equation}

\paragraph{Long-term feasible region.}
Because excess demand is either not admitted (TCP) or dropped (UDP), the queues
remain bounded even when the total demand exceeds the network capacity. The
long-term average delivered throughputs $\bar x_u$ therefore always satisfy
\begin{align}
  \bar x_u &\le \bar\lambda_u, && \forall u \in \Uset, \label{eq:demand} \\
  \sum_{u \in \Uset_v} \bar x_u &\le \frac{1}{T_{\mathrm{slot}}}\, \E\bigl[\, b_{\parent(v),v}(t)\, \rho_v(t) \,\bigr], && \forall (\parent(v),v) \in \Eset, \label{eq:link-cap} \\
  \sum_{v \in \Cset_n} \E\bigl[\, b_{n,v}(t) \,\bigr] &\le \E\bigl[\, \tilde B_n(t) \,\bigr], && \forall n \in \{0\} \cup \Nset, \label{eq:node-cap}
\end{align}
where the expectations are time averages over all slots, with $b_{n,v}(t) = 0$ and
$\tilde B_n(t) = 0$ outside $\Tset^{\mathrm{DL}}$. If the demand vector lies inside
this region, there exists a scheduling policy capable of supporting it; a
particular policy, for instance one with overly restrictive caps, may still fail to
do so. If the demand lies outside the region, no policy can support it, and the
policy determines which part of it is delivered.

\section{Optimization Problem}
\label{subsec:problem}

A policy $\pi$ selects in every epoch the cap tiers of all nodes,
$\bm a(k) = \{a_{n,v}(k)\}_{n \in \{0\}\cup\Nset,\, v \in \Cset_n}$. Over an
evaluation horizon of $K$ epochs, the system objective is the average sum
throughput:
\begin{subequations}
\label{eq:P1}
\begin{align}
  \textbf{(P1)}\quad \max_{\pi} \quad & \frac{1}{K} \sum_{k=0}^{K-1} \E_\pi\Bigl[ \sum_{u \in \Uset} x_u(k) \Bigr] \label{eq:P1-obj} \\
  \text{s.t.} \quad
  & a_{n,v}(k) \in \Aset, \quad \forall n, v \in \Cset_n, k, \label{eq:P1-action} \\
  & \text{\eqref{eq:du-budget}--\eqref{eq:cap}} \quad \text{(resources and scheduling)}, \\
  & \text{\eqref{eq:flow-conservation}--\eqref{eq:ue-tput}} \quad \text{(sources, queues and flow conservation)}, \\
  & \textstyle \frac{1}{K}\sum_{k} J\bigl(\bm\sigma(k)\bigr) \ge J_{\min} \quad \text{(optional fairness constraint)}. \label{eq:P1-fair}
\end{align}
\end{subequations}
The demand-satisfaction ratio and Jain's fairness index~\cite{jain1984fairness} are
\begin{equation}
  \sigma_u(k) = \min\Bigl\{ 1,\ \frac{x_u(k)}{\lambda_u(k)} \Bigr\}, \qquad
  J(\bm y) = \frac{\bigl(\sum_{i} y_i\bigr)^2}{|\bm y| \sum_i y_i^2},
  \label{eq:metrics}
\end{equation}
where $\bm\sigma(k)$ contains only the UEs with $\lambda_u(k) > 0$, and
$J(\bm y) = 1$ by convention when $\bm y = \bm 0$.

(P1) is the ideal centralized problem in which every DU, including DU$_0$, is
controlled. The schemes of this work keep the donor on PF and therefore solve the
restricted problem
\begin{equation}
  \textbf{(P1$'$)}\quad \text{(P1) with the additional constraint } a_{0,v}(k) = 1 \quad \forall v \in \Cset_0,\ \forall k .
  \label{eq:P1r}
\end{equation}

\paragraph{Control schemes.}
The schemes compared in this work are nested. Stage~1 is PF, i.e., $a_{n,v}(k) = 1$
for all $n$ and $v$. Stages~2--4 control the IAB nodes with the MARL formulation of
Section~\ref{subsec:marl}, use identical local learning, and differ only in the
federated aggregation of Section~\ref{subsec:fl}; they do not
enforce~\eqref{eq:P1-fair}. Stage~5 keeps the aggregation of Stage~4 and enforces a local surrogate of~\eqref{eq:P1-fair} by
Lagrangian relaxation: each agent uses the reward
\begin{equation}
  r_n^{\mathrm{L}}(k) = r_n(k) + \lambda_n \bigl( J_n(k) - J_{\min} \bigr),
  \label{eq:lagrangian}
\end{equation}
where $r_n(k)$ is the reward~\eqref{eq:reward} of Section~\ref{subsec:marl} and
$J_n(k)$ is Jain's index of the normalized per-child throughputs
$\min\{X_{n,v}(k)/X_{\max}, 1\}$, $v \in \Cset_n^{\act}(k)$. The multiplier is
updated by projected dual ascent after each training batch,
\begin{equation}
  \lambda_n \leftarrow \min\Bigl\{ \lambda_{\max},\ \bigl[ \lambda_n + \eta_\lambda \bigl( J_{\min} - \bar J_n \bigr) \bigr]^{+} \Bigr\},
  \label{eq:dual}
\end{equation}
where $\bar J_n$ is the batch average of $J_n(k)$ and $\lambda_n$ is initialized
to zero.

\paragraph{Why (P1) is not solved directly.}
First, the joint action space has size $\prod_{n} L^{|\Cset_n|}$, which is
exponential in the number of nodes and children. Second, actions are coupled across
hops and across time: the action of node $n$ changes the arrivals of
downstream nodes through~\eqref{eq:flow-conservation}, and queues carry these
effects into later epochs. Third, the statistics of the
channels and of the demands are unknown and non-stationary. Fourth, at the
near-real-time scale each DU observes only its own children. This work therefore
adopts a decentralized, learning-based approximation in which each node is an agent
acting on local observations and knowledge is shared through federated learning.

\section{Multi-Agent Reinforcement Learning Formulation}
\label{subsec:marl}

\paragraph{Game formulation.}
Every IAB node $n \in \Nset$ is an agent; the donor keeps PF. The interaction is
modelled as a partially observable stochastic game
$\langle \Nset, \mathcal{S}, \{\mathcal{O}_n\}, \{\mathcal{A}_n\}, \mathcal{P}, \{r_n\} \rangle$,
where the global state $s(k) \in \mathcal{S}$ contains all queues, channels and
occupancies, agent $n$ receives an observation $o_n(k) \in \mathcal{O}_n$, selects
$a_n(k) \in \mathcal{A}_n$, and receives a local reward $r_n(k)$. The system goal is
the team objective (P1), but the agents are trained independently with the local
surrogate rewards defined below. The extent to which the surrogate improves (P1) is
evaluated empirically.

\paragraph{Timeline.}
At time $t_k$, the start of epoch $k$, agent $n$ forms $o_n(k)$ from measurements
collected during epoch $k-1$ and from queue lengths sampled at $t_k$. It selects
$a_n(k)$, which is applied as the caps~\eqref{eq:cap} throughout epoch $k$. The
reward $r_n(k)$ is computed from measurements collected during epoch $k$ and becomes
available at $t_{k+1}$, together with $o_n(k+1)$.

\paragraph{Observation.}
Let $\Cset_n^{\act}(k) \subseteq \Cset_n$ be the children connected at $t_k$, and let
\begin{equation}
  X_{n,v}(k) = \sum_{t \in \Tset_k} S_{n,v}(t)
\end{equation}
be the bits sent by DU$_n$ to child $v$ during epoch $k$. The observation consists of
a set of per-child features and a node-level vector,
\begin{equation}
  o_n(k) = \Bigl( \bigl\{ \bm z_{n,v}(k) \bigr\}_{v \in \Cset_n^{\act}(k)},\ \bm g_n(k) \Bigr),
\end{equation}
with
\begin{equation}
  \bm z_{n,v}(k) = \Bigl( \tfrac{\log(1 + X_{n,v}(k-1))}{\log(1 + X_{\max})},\ \tfrac{q_v(t_k)}{q_{\max}},\ \tfrac{\log(1 + Q_{n,v}(t_k))}{\log(1 + Q_{\max})} \Bigr),
\end{equation}
i.e., the bits sent in the previous epoch (a throughput proxy), the most recently
reported MCS (a channel-quality proxy) and the RLC queue length (a demand proxy).
The node-level vector is
\begin{equation}
  \bm g_n(k) = \bigl( |\Cset_n^{\act}(k)| / C_{\max},\ \tilde\phi_n(k),\ \beta_n(t_k),\ \hat p_n(k),\ \hat c_n(k) \bigr),
\end{equation}
which contains the fraction of active children, the global fairness bias
$\tilde\phi_n$ defined below, the available PRB fraction $\beta_n$ of~\eqref{eq:beta}, and two relational
features,
\begin{align}
  \hat p_n(k) &= \clip\bigl( 0.5 + \kappa_p\,(\beta_{\parent(n)}(t_k) - \bar\beta_{\parent(n)}(t_k)),\ 0,\ 1 \bigr), \label{eq:phat}\\
  \hat c_n(k) &= \frac{\log\bigl(1 + \sum_{m \in \Cset_n^{\mathrm{I}}} \sum_{v \in \Cset_m} Q_{m,v}(t_k)\bigr)}{\log\bigl(1 + |\Cset_n^{\mathrm{I}}|\, Q_{\max}\bigr)}. \label{eq:chat}
\end{align}
Here $\hat p_n$ is the deviation of the parent's available fraction from its
exponentially smoothed value $\bar\beta_{\parent(n)}$, an early indicator of upstream
congestion, and $\hat c_n$ is the total queue of the node's IAB children, i.e., the
downstream demand. By convention $\hat p_n = 0.5$ when $\parent(n) = 0$ or when the
parent's value is unavailable, and $\hat c_n = 0$ when $\Cset_n^{\mathrm{I}} = \emptyset$
or when the children's values are unavailable. Both features require only the
exchange of scalars between neighbours in the tree. Under
Assumption~\ref{ass:orthogonal}, $\beta_n \equiv 1$ and hence $\hat p_n \equiv 0.5$;
these two features are constant on the testbed and become informative only for
in-band configurations, in which the available resources vary.

The Global xApp computes the fairness bias by comparing each node with the nodes of
the same role,
\begin{equation}
  \phi_n = \clip\Bigl( \frac{\bar X_{\role(n)}}{\bar X_n + \epsilon_0},\ \phi_{\min},\ \phi_{\max} \Bigr), \qquad
  \tilde\phi_n = \frac{\phi_n - \phi_{\min}}{\phi_{\max} - \phi_{\min}},
  \label{eq:fairness-bias}
\end{equation}
where $\bar X_n$ is the recent average throughput of node $n$ and
$\bar X_{\role(n)}$ is the average over the nodes of the same role. The comparison is
restricted to the same role because the throughput of a relay includes the backhaul
traffic it forwards and is therefore of a different magnitude.

\paragraph{Action.}
If $|\Cset_n^{\act}(k)| \ge 2$, the action is
$a_n(k) = \bigl( a_{n,v}(k) \bigr)_{v \in \Cset_n^{\act}(k)} \in \Aset^{|\Cset_n^{\act}(k)|}$,
converted to caps by~\eqref{eq:cap}. If $|\Cset_n^{\act}(k)| \le 1$ there is no
contention between children, a cap could only leave resources unused, and the node
applies $a_{n,v}(k) = 1$ without querying the policy.

\paragraph{Reward.}
For $\Cset_n^{\act}(k) \neq \emptyset$,
\begin{equation}
  r_n(k) = \frac{1}{|\Cset_n^{\act}(k)|} \sum_{v \in \Cset_n^{\act}(k)} \min\Bigl\{ \frac{X_{n,v}(k)}{X_{\max}},\ 1 \Bigr\},
  \label{eq:reward}
\end{equation}
i.e., the normalized throughput that DU$_n$ delivers to its children during the
epoch in which the action is applied; epochs without active children are not used
for training. This is a surrogate for (P1) and is not guaranteed to be aligned with
it. For a relay, sending more backhaul traffic to a child increases the reward but
does not necessarily increase the end-to-end throughput, because the traffic may
only move into a downstream queue. In addition,
the truncation and averaging in~\eqref{eq:reward} differ from the sum
in~\eqref{eq:P1-obj}.

\paragraph{Policy and value architecture.}
Because $|\Cset_n^{\act}|$ differs across nodes and epochs and children have no
natural order, the policy factorizes over children with a shared per-child network,
\begin{align}
  \pi_\theta\bigl(a_n \mid o_n\bigr) &= \prod_{v \in \Cset_n^{\act}} \pi_\theta\bigl( a_{n,v} \mid \bm\psi_{n,v} \bigr), \label{eq:actor}\\
  \bm\psi_{n,v} &= \Bigl[ \bm z_{n,v},\ \bm g_n,\ \operatorname*{mean}_{v' \in \Cset_n^{\act} \setminus \{v\}} \bm z_{n,v'},\ \max_{v' \in \Cset_n^{\act} \setminus \{v\}} \bm z_{n,v'} \Bigr], \notag
\end{align}
where $\pi_\theta(\cdot \mid \bm\psi)$ is a categorical distribution over the $L$
tiers and the element-wise leave-one-out pooling is over a non-empty set since
$|\Cset_n^{\act}| \ge 2$. This form is equivariant to permutations of the children.
The critic is permutation invariant,
\begin{equation}
  V_\vartheta(o_n) = \varrho_\vartheta\Bigl( \operatorname*{mean}_{v \in \Cset_n^{\act}} \varphi_\vartheta(\bm z_{n,v}),\ \max_{v \in \Cset_n^{\act}} \varphi_\vartheta(\bm z_{n,v}),\ \bm g_n \Bigr).
  \label{eq:critic}
\end{equation}
With this structure, parameters of the same dimension apply to a node at any
position in the tree with any number of children, which is required for averaging
parameters across all nodes in Section~\ref{subsec:fl}.

\paragraph{Initialization.}
The policy is initialized so that $\pi_\theta(\alpha_L \mid \bm\psi) = p_0$ for every
input, so the initial policy coincides with PF on each child with probability $p_0$,
and learning acts as a residual correction to PF. The critic is pretrained by
regression on $(o_n(k), r_n(k))$ samples collected while all nodes run PF.

\paragraph{Learning objective.}
This work uses a myopic objective: each agent maximizes the expected immediate
reward $\E[r_n(k)]$, which corresponds to a discount factor of zero. This is a
contextual-bandit-style approximation. It neglects the effects of an action on
later epochs through queues and downstream arrivals, which
are present in the system; the approximation is motivated by the fact that the cap
of an epoch mainly affects the bits sent in that epoch. Let
\begin{equation}
  \xi_n(k) = \Ind\Bigl[ |\Cset_n^{\act}(k)| \ge 2 \ \wedge\ \max_{v \in \Cset_n^{\act}(k)} Q_{n,v}(t_k) \ge Q_{\mathrm{th}} \Bigr]
  \label{eq:contended}
\end{equation}
indicate contention at decision time. A cap can reduce the throughput of a child
only if the child needs more resources than the cap allows, which is most likely
when several children are active and at least one of them has a large backlog; the
indicator is a heuristic proxy for this condition and is evaluated before the
action, so that it does not depend on the action. Let $\mathcal{K}_n$ be the set of pairs $(k,v)$ in a
training batch with $\xi_n(k) = 1$ and $v \in \Cset_n^{\act}(k)$. The actor
minimizes $-\mathcal{L}_n(\theta)$ with
\begin{equation}
  \mathcal{L}_n(\theta) = \frac{1}{|\mathcal{K}_n|} \sum_{(k,v) \in \mathcal{K}_n} \min\Bigl( \eta_{n,v}(k)\, \hat A_n(k),\ \clip\bigl(\eta_{n,v}(k), 1 - \epsilon_{\mathrm{clip}}, 1 + \epsilon_{\mathrm{clip}}\bigr)\, \hat A_n(k) \Bigr) + c_H\, \bar{\mathcal{H}}_n(\theta),
  \label{eq:ppo}
\end{equation}
where
$\eta_{n,v}(k) = \pi_\theta(a_{n,v}(k) \mid \bm\psi_{n,v}(k)) / \pi_{\theta_{\mathrm{old}}}(a_{n,v}(k) \mid \bm\psi_{n,v}(k))$
is the per-child probability ratio with respect to the behaviour policy that
generated the sample,
$\hat A_n(k) = \clip\bigl( (r_n(k) - V_\vartheta(o_n(k)))/\hat\sigma_r,\ -3,\ 3 \bigr)$
with $\hat\sigma_r$ a running standard deviation of the reward, and
$\bar{\mathcal{H}}_n$ is the average per-child policy entropy. Uncontended samples
are used only to train the critic. Since the joint policy~\eqref{eq:actor} is a
product over children, the joint-action ratio of standard PPO would be
$\prod_{v} \eta_{n,v}(k)$; \eqref{eq:ppo} instead clips each factor separately and
shares the node-level advantage among the children, i.e., it is a factor-wise
surrogate of the PPO objective~\cite{schulman2017ppo}.

\section{Hierarchical Federated Learning}
\label{subsec:fl}

Each agent holds local parameters $\Theta_n = (\theta_n, \vartheta_n)$ and updates
them on local data. Every $\Delta_{\mathrm{FL}}$ seconds a federated round $r$ takes
place. Node $n$ reports $\Theta_n^{(r)}$ and the weight $w_n^{(r)} \ge 0$, the number
of contended samples in its current training data that can update the actor. Let
$\tilde\Theta_n^{(r-1)}$ be the parameters held by node $n$ at the end of round
$r-1$, and $\Delta_n^{(r)} = \Theta_n^{(r)} - \tilde\Theta_n^{(r-1)}$ its local
update. If $W^{(r)} = \sum_{n \in \Nset} w_n^{(r)} = 0$, the round is skipped and
all parameters are left unchanged. Otherwise, one of the following aggregation
rules is applied. Unless stated otherwise, every node, including those with
$w_n^{(r)} = 0$, receives the broadcast.

\paragraph{Stage 2: flat FedAvg.}
The baseline is federated averaging (FedAvg)~\cite{mcmahan2017fedavg},
\begin{equation}
  \Theta_{\mathrm{glob}}^{(r)} = \sum_{n \in \Nset} \frac{w_n^{(r)}}{W^{(r)}}\, \Theta_n^{(r)}, \qquad \Theta_n \leftarrow \Theta_{\mathrm{glob}}^{(r)} \quad \forall n \in \Nset.
  \label{eq:fedavg}
\end{equation}

\paragraph{Stage 3: topology-hierarchical FedAvg with Hedge.}
The branches~\eqref{eq:branch} form the first level of aggregation.
For branch $b$, let $\Nset_b^{+} = \{n \in \Nset_b : w_n^{(r)} > 0\}$ be its
contributing members.

\emph{Divergence fallback.} Following the idea of clustered federated
learning~\cite{sattler2021cfl}, if $|\Nset_b^{+}| \ge 3$, the alignment of each member
with the rest of its branch is
\begin{equation}
  \operatorname{sim}_n^{(r)} = \cos\Bigl( \Delta_n^{(r)},\ \sum_{n' \in \Nset_b^{+} \setminus \{n\}} w_{n'}^{(r)} \Delta_{n'}^{(r)} \Bigr), \qquad n \in \Nset_b^{+},
  \label{eq:cfl}
\end{equation}
with $\cos(\cdot,\cdot) = 1$ by convention if either argument is the zero vector.
The excluded set is $\mathcal{X}_b^{(r)} = \{ n \in \Nset_b^{+} : \operatorname{sim}_n^{(r)} < \tau_{\mathrm{CFL}} \}$
if this set is a proper subset of $\Nset_b^{+}$, and $\mathcal{X}_b^{(r)} = \emptyset$
otherwise or if $|\Nset_b^{+}| < 3$.

\emph{Main groups.} The main group of branch $b$ is
$G_b = \Nset_b^{+} \setminus \mathcal{X}_b^{(r)}$, which is non-empty whenever
$\Nset_b^{+} \neq \emptyset$. Let $\Bset^{(r)} = \{ b \in \Bset : \Nset_b^{+} \neq \emptyset \}$.
For $b \in \Bset^{(r)}$ the weight and the within-branch average are
\begin{equation}
  W_{G_b}^{(r)} = \sum_{n \in G_b} w_n^{(r)} > 0, \qquad
  \bar\Theta_{G_b}^{(r)} = \sum_{n \in G_b} \frac{w_n^{(r)}}{W_{G_b}^{(r)}}\, \Theta_n^{(r)} .
\end{equation}

\emph{Hedge weights.} Following the Hedge algorithm~\cite{freund1997hedge}, each
branch keeps a weight $\zeta_b > 0$, initialized to $\zeta_b^{(0)} = 1$. For every
$b \in \Bset^{(r)}$,
\begin{equation}
  \zeta_b^{(r)} = \zeta_b^{(r-1)}\, e^{-\eta_H\, \ell_b^{(r)}},
\end{equation}
where $\ell_b^{(r)}$ is the mean squared error of the critic of
$\bar\Theta_{G_b}^{(r)}$ on the contended samples of the members of $G_b$; the
weights of the other branches are unchanged. The global parameters are
\begin{equation}
  \hat\zeta_b^{(r)} = \frac{W_{G_b}^{(r)}\, \zeta_b^{(r)}}{\sum_{b' \in \Bset^{(r)}} W_{G_{b'}}^{(r)}\, \zeta_{b'}^{(r)}}, \qquad
  \Theta_{\mathrm{glob}}^{(r)} = \sum_{b \in \Bset^{(r)}} \hat\zeta_b^{(r)}\, \bar\Theta_{G_b}^{(r)} .
  \label{eq:hedge}
\end{equation}

\emph{Broadcast.} Let $\mathcal{X}^{(r)} = \bigcup_{b} \mathcal{X}_b^{(r)}$. The nodes
are updated as
\begin{equation}
  \Theta_n \leftarrow
  \begin{cases}
    \Theta_n^{(r)}, & n \in \mathcal{X}^{(r)}, \\
    \Theta_{\mathrm{glob}}^{(r)}, & n \in \Nset \setminus \mathcal{X}^{(r)} .
  \end{cases}
  \label{eq:stage3-broadcast}
\end{equation}
An excluded node thus forms a singleton cluster for this round: its update does not
enter $\Theta_{\mathrm{glob}}^{(r)}$, and it keeps its own parameters instead of
being overwritten. The exclusion is re-evaluated in every round, so a node rejoins
the global model as soon as its update is aligned with its branch again. Stage~3 is
therefore the first scheme in which the parameters of different nodes may differ.
For deeper trees the within-branch average can be applied recursively along
subtrees; this work uses two levels.

\paragraph{Nesting of the stages.}
With $\tau_{\mathrm{CFL}} = -1$ no node is excluded, since
$\operatorname{sim}_n^{(r)} \ge -1$, so every node receives
$\Theta_{\mathrm{glob}}^{(r)}$. With $\eta_H = 0$ the Hedge weights remain equal.
Then
\begin{equation}
  \Theta_{\mathrm{glob}}^{(r)} = \sum_{b \in \Bset^{(r)}} \frac{W_{G_b}^{(r)}}{W^{(r)}} \sum_{n \in G_b} \frac{w_n^{(r)}}{W_{G_b}^{(r)}}\, \Theta_n^{(r)} = \sum_{n \in \Nset} \frac{w_n^{(r)}}{W^{(r)}}\, \Theta_n^{(r)},
\end{equation}
which is Stage 2, because nodes with $w_n^{(r)} = 0$ contribute to neither side.

\section{Mapping to the O-RAN Architecture}
\label{subsec:oran}

Table~\ref{tab:oran-mapping} maps the model components to their deployment. The
testbed implements the E2 interface and a Near-RT RIC, but no Non-RT RIC, A1
interface or service management and orchestration framework. The xApps are
standard Near-RT RIC applications. The components called ``Local rApp'' and
``Global rApp'' in this work are custom services and not rApps in the O-RAN sense,
which reside in the Non-RT RIC. The Local rApp performs policy inference for the
Local xApp over inter-process communication and trains the local model in the
background; the Global rApp is the FL server. Both are co-located with the
Near-RT RIC. Functionally, model training and federated aggregation correspond to
non-real-time tasks. The periods in the table are update periods of the
applications, not claims about the control-loop time scales defined by O-RAN.

\begin{table}[htbp]
\centering
\small
\caption{Mapping of model components to their deployment and update periods.}
\label{tab:oran-mapping}
\begin{tabular}{@{}p{0.30\textwidth}p{0.38\textwidth}p{0.24\textwidth}@{}}
\toprule
Model component & Deployment & Period on the testbed \\
\midrule
PF scheduling within $\tilde B_n(t)$ & O-DU MAC (E2 node) & slot, 0.5 ms \\
Reporting of per-child statistics & E2SM-MAC indication & 100 ms \\
Observation $o_n(k)$ and cap enforcement & Local xApp (Near-RT RIC) & epoch $\Delta$ = 1 s \\
Inference of $\pi_\theta$ & Local rApp (custom service) & epoch $\Delta$ = 1 s \\
Local update~\eqref{eq:ppo} & Local rApp (custom service) & 60 s \\
Fairness bias $\tilde\phi_n$ & Global xApp (Near-RT RIC) & 2 s \\
Federated aggregation~\eqref{eq:fedavg}--\eqref{eq:stage3-broadcast} & Global rApp (custom FL server) & $\Delta_{\mathrm{FL}}$ = 180 s \\
\bottomrule
\end{tabular}
\end{table}

\paragraph{Testbed parameters.}
The testbed uses $\mu = 1$ ($T_{\mathrm{slot}} = 0.5$~ms), $B_n = 106$, $\nu = 1$,
$Q_m^{\max} = 8$, $R^{\max} = 948/1024$, $f = 1$ and $\mathrm{OH} = 0.14$,
for which~\eqref{eq:peak} gives $r^{\max} \approx 227$~Mbps. The TDD pattern has a
period of $P = 10$ slots with 7 downlink slots, one special slot with 6 downlink and
4 uplink symbols, and 2 uplink slots. Hence
$\delta^{\mathrm{DL}}_{\mathrm{symb}} = 104/140$ and the downlink peak is about
$169$~Mbps (about $159$~Mbps if only full downlink slots carry data). All DUs,
including those serving backhaul links, use $B_n = 106$, and $\beta_n \equiv 1$. The cap tiers are $\Aset = \{0.3, 0.5, 0.7, 0.85, 1.0\}$, with
$p_0 = 0.8$, $b_{\min} = 5$, $Q_{\mathrm{th}} = 10^5$ bytes and
$\epsilon_{\mathrm{clip}} = 0.2$.

\section{Notation}
\label{subsec:notation}

\begin{longtable}{@{}p{0.26\textwidth}p{0.68\textwidth}@{}}
\toprule
Symbol & Meaning \\
\midrule
\endhead
$0$, $\Nset$, $\Uset$ & Donor; IAB nodes $\{n_1,\dots,n_N\}$; UEs $\{u_1,\dots,u_U\}$ \\
$\parent(v)$, $\Cset_n$ & Parent; children $\Cset_n = \Cset_n^{\mathrm{I}} \cup \Cset_n^{\mathrm{U}}$ \\
$h(v)$, $H$ & Depth (hop count); maximum depth \\
$\Uset_v$, $\Nset_v$, $\mathcal{P}(u)$ & UEs and IAB nodes in the subtree of $v$; downlink path of UE $u$ \\
$\role(n)$, $\br(n)$, $\Bset$ & Role of a node; its branch; set of branches \\
$t$, $T_{\mathrm{slot}}$, $k$, $\Tset_k$, $\Delta$ & Slot; slot length; epoch; slots of epoch $k$; epoch length \\
$\delta^{\mathrm{DL}}$, $\delta^{\mathrm{UL}}$ & Fractions of downlink- and uplink-capable slots \\
$B_n$, $\tilde B_n(t)$ & Total PRBs of DU$_n$; PRBs available to DU$_n$ in slot $t$ \\
$\beta_n(t)$ & Available fraction of the PRBs of DU$_n$ ($\equiv 1$ in out-of-band mode) \\
$\gamma_v$, $q_v$, $\rho_v$ & SINR; MCS index; bits per PRB per downlink slot \\
$b_{n,v}(t)$, $\bar b_{n,v}(k)$ & PRBs allocated by DU$_n$ to child $v$; cap \\
$\Aset$, $a_{n,v}(k)$ & Set of cap tiers ($\alpha_L = 1$ is PF); selected tier \\
$\lambda_u$, $A_u$ & Target rate of UE $u$; bits injected at the donor \\
$Q_n^{(u)}$, $Q_{n,v}$, $S_{n,v}$, $D_{n,v}$ & Per-UE queue; aggregated queue towards child $v$; bits sent; bits dropped \\
$Q^{\mathrm{buf}}$ & Buffer capacity per child \\
$x_u(k)$, $\sigma_u(k)$, $J(\cdot)$ & Delivered throughput; demand-satisfaction ratio; Jain's index \\
$\bar c_{\parent(n),n}$, $\bar c_n$, $F_n$ & Backhaul and access capacities of node $n$; traffic through $n$ \\
$o_n$, $\bm z_{n,v}$, $\bm g_n$ & Observation; per-child feature; node-level feature \\
$\hat p_n$, $\hat c_n$, $\tilde\phi_n$ & Parent trend; downstream demand; fairness bias \\
$r_n$, $\xi_n$, $X_{n,v}$ & Local reward; contention indicator; bits sent per epoch \\
$\theta_n$, $\vartheta_n$ & Actor and critic parameters \\
$w_n^{(r)}$, $\Delta_n^{(r)}$ & FL weight; local update \\
$G_b$, $\mathcal{X}_b^{(r)}$, $\tau_{\mathrm{CFL}}$ & Main group of branch $b$; excluded nodes; divergence threshold \\
$\zeta_b$, $\eta_H$ & Hedge weight; Hedge learning rate \\
$\lambda_n$, $J_{\min}$ & Lagrange multiplier of Stage 5; fairness target \\
\bottomrule
\end{longtable}

\end{EnChapter}
